\documentclass[11pt]{article}
\usepackage{amsmath,amssymb,amsthm}
\usepackage[margin=1in]{geometry}
\theoremstyle{plain}
\newtheorem{theorem}{Theorem}
\newtheorem{lemma}[theorem]{Lemma}
\newtheorem{corollary}[theorem]{Corollary}
\newtheorem{proposition}[theorem]{Proposition}
\theoremstyle{definition}
\newtheorem{definition}[theorem]{Definition}
\theoremstyle{remark}
\newtheorem{remark}[theorem]{Remark}
\newcommand{\SYT}{\operatorname{SYT}}
\newcommand{\cont}{\operatorname{cont}}
\newcommand{\Des}{\operatorname{Des}}
\newcommand{\dD}{\Delta D}
\title{Descent transport and the Tightness Lemma:\\
a reformulation of the Core Lemma residual}
\author{Clio}
\date{2026-05-22}
\begin{document}
\maketitle

\begin{abstract}
Working from the arch-telescoping reduction of the Core Lemma
($\mathsf{mincost}(\lambda)\ge n(\lambda)$), I give a transparent reformulation of the
per-arch sign data: $\sigma$ is the change in descent number, and a $\sigma=0$ swap
\emph{transports a single descent by one position}, with $e=\alpha-\beta$ recording the
direction. Using this language I prove that the Forest Inequality's downward induction
\emph{closes at every non-root arch} using only the proved per-arch inequality, provided
internal arches satisfy $\dD<0$. That hypothesis is exactly the content of a single clean
\textbf{Tightness Lemma} (a non-root arch with $\dD>0$ has width $1$, hence is a leaf),
which also subsumes the rule (S2). The only residual case is a $\dD>0$ \emph{root}. I show
by explicit counterexample that this residual is \emph{not} closed by the rules
(R),(S2) and the realized $\sigma$-transition bans alone: it provably requires the
\textbf{comparator-step coupling} $|i_c-i_p|\le 2$, an ingredient absent from the previous
reduction. This corrects and sharpens the open-problem list.
\end{abstract}

\section{Setup and notation}
We use the notation of the arch-telescope note. $\lambda\vdash n$, $\cont(r,c)=c-r$,
$\cont_T(v)$ the content of the cell holding label $v$ in $T\in\SYT(\lambda)$, and
$\delta_T(i)=[\cont_T(i{+}1)<\cont_T(i)]$ for $1\le i\le n-1$. The \emph{descent set} is
$\Des(T)=\{i:\delta_T(i)=1\}$. A closed feasible walk on the staircase word
$\mathbf w_0=(1)(2\,1)\cdots(n{-}1\cdots1)$ produces, at each comparator $i$, a STAY or a
legal SWAP $U\mapsto s_iU$ (exchanging the cells of labels $i,i{+}1$), legal iff
$|d_U(i)|\ge2$ with $d_U(i)=\cont_U(i{+}1)-\cont_U(i)$; feasibility forbids any visited
before-state with $d=-1$.

The swaps match into nested \emph{arches}; the nesting forest is a union of chains. For an
arch $a$ at comparator $i$ with boundary $U$ (just before its opening swap) and inner
$V=s_iU$, set $x_j=\cont_U(j)$ and
\[
\gamma=\delta_V(i)-\delta_U(i),\quad
\alpha=\delta_V(i{-}1)-\delta_U(i{-}1),\quad
\beta=\delta_V(i{+}1)-\delta_U(i{+}1),
\]
(with $\alpha=0$ if $i=1$, $\beta=0$ if $i{+}1=n$), and $\sigma=\gamma+\alpha+\beta$,
$e=\alpha-\beta$, $ni=n-i$, $w=b_2-b_1$ (closing minus opening block). The proved
\emph{sign rule} gives $\gamma=1-2\delta_U(i)=\pm1$, $\alpha,\beta\in\{0,-\gamma\}$,
$\sigma\in\{-1,0,1\}$; $\sigma=\pm1\Rightarrow\alpha=\beta$ and $\sigma=0\Rightarrow
|e|=1$. The proved per-arch data: $\Delta_a=w\sigma$, $\dD_a=ni\,\sigma+e$,
$1\le w\le ni-1$, and the per-arch inequality
\begin{equation}\label{eq:perarch}
\dD_a\ne0;\qquad \dD_a<0\Rightarrow\Delta_a\ge \dD_a+1;\qquad \dD_a>0\Rightarrow\Delta_a\ge0.
\end{equation}
The \emph{boundary identity} along a chain is $U_{c}=V_{p}=s_{i_p}U_p$, and the proved
\emph{width-coupling} lemma gives $w_{c}\le w_{p}-1$ (so widths strictly decrease down a
chain).

\section{Reformulation: $\sigma$ is descent transport}

\begin{theorem}[Descent-transport]\label{thm:transport}
For any legal swap $U\mapsto V=s_iU$:
\begin{enumerate}
\item $\sigma=|\Des(V)|-|\Des(U)|$. Thus $\sigma=+1$ creates a descent, $\sigma=-1$
destroys one, $\sigma=0$ preserves the descent number.
\item If $\sigma=0$ then exactly one descent moves by one position: there are positions
$p\ne p'$ with $\{p,p'\}\in\{\{i{-}1,i\},\{i,i{+}1\}\}$ such that
$\Des(V)=\big(\Des(U)\setminus\{p\}\big)\cup\{p'\}$, and
\[
e \;=\; \operatorname{sign}(p-p')\;=\;-\operatorname{sign}(p'-p).
\]
That is, $e=+1$ iff the descent moves left ($p'<p$) and $e=-1$ iff it moves right.
\end{enumerate}
\end{theorem}
\begin{proof}
The swap changes only the contents of labels $i,i{+}1$, hence only the descent indicators
$\delta(\cdot)$ at the comparators $i{-}1,i,i{+}1$ can change; their total change is
$\alpha+\gamma+\beta=\sigma$, proving (1).

For (2), assume $\sigma=0$. By the sign rule $\gamma=\pm1$ and exactly one of
$\alpha,\beta$ equals $-\gamma$, the other $0$. Each nonzero indicator-change is $\pm1$, so
two adjacent comparators toggle in opposite senses: one descent is removed and one created,
at positions differing by $1$ and both incident to $i$. Concretely:
\begin{itemize}
\item $\gamma=+1,\alpha=-1$: $\delta(i)$ rises ($p'=i$ created), $\delta(i{-}1)$ falls
($p=i{-}1$ removed); $p\to p'$ is $i{-}1\to i$ (right), $e=\alpha-\beta=-1$.
\item $\gamma=-1,\alpha=+1$: $i\to i{-}1$ (left), $e=+1$.
\item $\gamma=+1,\beta=-1$: $i{+}1\to i$ (left), $e=+1$.
\item $\gamma=-1,\beta=+1$: $i\to i{+}1$ (right), $e=-1$.
\end{itemize}
In every case $e=\operatorname{sign}(p-p')$, as claimed.
\end{proof}

\begin{corollary}\label{cor:dDsign}
$\dD_a>0$ iff $\sigma=+1$ (a descent is created) or $\sigma=0,e=+1$ (a descent moves left);
$\dD_a<0$ iff $\sigma=-1$ or $\sigma=0,e=-1$. Indeed $\dD=ni\,\sigma+e$ with $ni\ge2$ and
$e=0$ when $\sigma\ne0$, $|e|=1$ when $\sigma=0$.
\end{corollary}

\section{The Forest induction closes off the root}

Recall the Forest quantities for a single chain $a_1$ (root) $,\dots,a_m$ (leaf):
$T=\sum_j\Delta_{a_j}$, $M_m=\dD_{a_m}$, $M_j=\dD_{a_j}+\min(0,M_{j+1})$, and the target
\[
\textbf{Forest Inequality (single chain):}\qquad T\ \ge\ \mu:=\min(0,M_1).
\]
Write $R_i=\sum_{j\ge i}\Delta_{a_j}$ and $S_i=\min(0,M_i)$ (so $R_1=T$, $S_1=\mu$).

\begin{lemma}[Non-root steps close unconditionally]\label{lem:induct}
Suppose every \emph{internal} arch $a_i$ (i.e.\ $1<i<m$) has $\dD_{a_i}<0$. Then
$R_i\ge S_i$ for all $i\ge2$, and moreover the inductive step is valid at every $i\ge2$
using only \eqref{eq:perarch}.
\end{lemma}
\begin{proof}
Downward induction on $i$. \emph{Base} $i=m$: $R_m=\Delta_{a_m}$, $S_m=\min(0,\dD_{a_m})$.
If $\dD_{a_m}>0$ then $\Delta_{a_m}\ge0=S_m$ by \eqref{eq:perarch}; if $\dD_{a_m}<0$ then
$\Delta_{a_m}\ge\dD_{a_m}+1>\dD_{a_m}=S_m$. \emph{Step} $2\le i<m$ (internal, so
$\dD_{a_i}<0$): since $S_{i+1}=\min(0,M_{i+1})\le0$ and $\dD_{a_i}<0$, we have
$\dD_{a_i}+S_{i+1}<0$, hence $S_i=\dD_{a_i}+S_{i+1}$. By the induction hypothesis
$R_{i+1}\ge S_{i+1}$, so
$R_i=\Delta_{a_i}+R_{i+1}\ge\Delta_{a_i}+S_{i+1}\ge(\dD_{a_i}+1)+S_{i+1}>S_i$,
using \eqref{eq:perarch}. \end{proof}

\begin{remark}
Lemma~\ref{lem:induct} isolates the entire difficulty of the Forest Inequality into
\emph{(i)} the hypothesis ``internal arches have $\dD<0$'' and \emph{(ii)} the final root
step $i=1$. The internal-$\dD<0$ hypothesis is supplied precisely by the Tightness Lemma
below; the root step is the genuine residual.
\end{remark}

\section{The Tightness Lemma (subsumes (S2))}

\begin{definition}
An arch is \emph{tight} if $w=1$. By width-coupling, widths strictly decrease down a chain,
so a tight non-root arch must be a leaf.
\end{definition}

\begin{lemma}[Tightness; \emph{verified $n\le7$, proof open}]\label{lem:tight}
A non-root arch $a$ with $\dD_a>0$ has $w_a=1$.
\end{lemma}

\noindent\emph{Status.} Verified with zero exceptions over all $27$ non-root arches with
$\dD>0$ across all feasible closed walks $n\le7$. Equivalently (Corollary~\ref{cor:dDsign}):
a non-root arch that \emph{creates} a descent ($\sigma=+1$) or \emph{moves a descent left}
($\sigma=0,e=+1$) is tight. The block identity
$w_a=w_p-(b_1^a-b_1^p)-(b_2^p-b_2^a)$ holds in general; the open content is why $\dD_a>0$
forces the offsets to sum to $w_p-1$.

\begin{corollary}[(S2)]\label{cor:s2}
A $(\sigma=0,e=+1)$ arch with a parent is a leaf; in particular no chain contains
$(\sigma=0)\to(\sigma=0,e=+1)\to(\cdot)$. Hence the named residual violator
$[(\sigma{=}0,e{=}{+}1),(\sigma{=}0,e{=}{+}1),(\sigma{=}-1,ni{=}2,w{=}1)]$ is unrealizable.
\end{corollary}
\begin{proof}
Such an arch has $\dD=+1>0$ and a parent (non-root), so by Lemma~\ref{lem:tight} $w=1$;
being non-root and tight it is a leaf (width strict decrease). The middle arch of the named
chain is a non-root $(\sigma=0,e=+1)$ with a child, contradiction.
\end{proof}

\begin{theorem}[Forest Inequality, modulo the root case]\label{thm:forestmod}
Assume Lemma~\ref{lem:tight}. Then for any single chain the Forest Inequality $T\ge\mu$ holds
\emph{unless} possibly the root has $\dD_{a_1}>0$.
\end{theorem}
\begin{proof}
By Lemma~\ref{lem:tight} a non-root arch with $\dD>0$ is tight, hence a leaf; so every
internal arch ($1<i<m$) has $\dD<0$. Lemma~\ref{lem:induct} gives $R_2\ge S_2$ and, if also
$\dD_{a_1}<0$, the same step at $i=1$ gives $R_1\ge S_1=\mu$. The only uncovered case is
$\dD_{a_1}>0$. (If $m=1$ and $\dD_{a_1}>0$ then $\mu=0$ and $T=\Delta_{a_1}\ge0$ by
\eqref{eq:perarch}, so even this is fine for $m=1$.)
\end{proof}

\section{The root residual needs the comparator-step coupling}

The remaining case is a chain whose root creates a descent or moves one left. Here the
previous rule-set is \emph{insufficient}.

\begin{proposition}[Insufficiency of (R),(S2),bans]\label{prop:counter}
There is an abstract chain satisfying width strict-decrease, the per-arch
inequality~\eqref{eq:perarch}, the Tightness filter, and \emph{all} of the realized
$\sigma$-transition bans \textnormal{(}no $-1\to-1$, no $+1\to+1$, and \textnormal{(R)}\textnormal{)},
yet violating the Forest Inequality:
\[
\big[(\sigma,e,ni,w)\big]=\big[(1,0,10,5),(-1,0,5,4),(0,-1,4,3),(-1,0,3,2),(0,-1,2,1)\big],
\]
with $T=-1<\mu=0$. It is excluded only because its comparator step
$i_c-i_p=ni_p-ni_c=10-5=5$ violates the realized bound $|i_c-i_p|\le2$.
\end{proposition}

\begin{remark}[Scope correction]\label{rmk:scope}
Proposition~\ref{prop:counter} shows the open list cannot be ``(R) and (S2)'' alone, nor
even those plus the $\sigma$-transition bans. By Theorem~\ref{thm:forestmod} the genuine
residual of the Core Lemma is:
\begin{itemize}
\item[\textbf{(T)}] the \textbf{Tightness Lemma} (Lemma~\ref{lem:tight}); and
\item[\textbf{(C)}] the \textbf{comparator-step coupling} $|i_c-i_p|\le2$ along a chain
(equivalently $ni_c\in\{ni_p+1,ni_p-1,ni_p-2\}$), which controls the $\dD>0$ root case;
\end{itemize}
together with the still-open Structural Lemma (arches nest into chains; $\le2$ roots). Both
(T) and (C) are read off the boundary identity $U_c=V_p=s_{i_p}U_p$ and are verified with
zero exceptions for $n\le7$. Notably (R) and (S2) are \emph{not} the right primitive
targets: (S2) is a corollary of (T), and (R) is not needed by the induction at all.

What is established here is that (C) is \emph{necessary} (Proposition~\ref{prop:counter})
and (T) is necessary (it supplies internal-$\dD<0$, without which the named violator and
the $56$ abstract violators survive). \emph{Sufficiency} of $\{$(T),(C)$\}$ for the
root-$\dD>0$ closure is not proved here; it must be checked that no further coupling is
needed. The real-walk Forest Inequality holds for all $n\le7$, and all real chains satisfy
(C); this is the evidence that $\{$(T),(C)$\}$ (with the Structural Lemma) is the right and
complete residual, but a closed-form proof of the root step from (C) remains open.
\end{remark}

\section{Status}
\textbf{Proved unconditionally:} Theorem~\ref{thm:transport} (descent transport),
Corollary~\ref{cor:dDsign}, Lemma~\ref{lem:induct} (non-root induction step),
Theorem~\ref{thm:forestmod} (Forest Inequality modulo Tightness and the root case),
Corollary~\ref{cor:s2} ((S2) from Tightness), and Proposition~\ref{prop:counter}.\\
\textbf{Verified $n\le7$, proof open:} Lemma~\ref{lem:tight} (Tightness) and the
comparator-step coupling (C); the Structural Lemma.\\
\textbf{Remaining gap:} (T) and (C); a closed-form proof that (C) closes the root-$\dD>0$
step (necessity is proved in Prop.~\ref{prop:counter}; sufficiency is verified $n\le7$ but
not proved); and the $\le2$-root slack absorption. With these the Core Lemma follows modulo
the Structural Lemma.
\end{document}
